\section{Methods}
\label{sec:methods}
% Every operational choice needed to reproduce the numbers, in the order the
% data flows. The development split is mentioned once; no development number
% is reported.

\subsection{Rule set}
\label{sec:rules}

The rules were taken from a single, widely read source: the ``IR Spectral Interpretation Workshop'' column by Brian C. Smith in \emph{Spectroscopy}, in its five ``Big Review'' summaries (hydrocarbons, the C-O bond, carbonyl compounds in two parts, organic nitrogen compounds)~\cite{smith2025bigreview4,smith2025bigreview5,smith2025bigreview6,smith2025bigreview7,smith2026bigreview8} and seven topical columns on benzene rings, alkenes, alkynes, aldehydes, nitriles, nitro groups and amine salts~\cite{smith2016benzene,smith2016alkenes,smith2017alkynes,smith2017aldehydes,smith2019nitriles,smith2020nitro,smith2019aminesalts}. We chose one source rather than a compilation because the study needs each rule to be exactly what one author wrote, with its window, its intensity word and its conditions, and because these columns are the reference many practising chemists use today. The tables were transcribed into a machine-readable form: one row per rule, with the group it refers to, the lower and upper wavenumber of the window, the intensity (``strong'', ``medium'' or unstated) and shape (``sharp'', ``broad'' or unstated) qualifiers as printed, and, for the benzene substitution patterns, the second window the rule requires or forbids. Where the source gives a centre and a tolerance the window is the centre plus and minus the tolerance; the one rule given as an approximate centre with no tolerance (secondary aromatic amine N-H, ``about 3400'') was assigned the 40~\wn{} width of its neighbours in the same table, and this is the only value not printed in the source.

The table holds 106 rules. Four refer to cis and trans alkenes and cannot be scored because most deposited structures carry no stereochemistry, and one (the overtone and combination bands of amine salts) covers three salt types and has no single structural target; these five are kept in the table as not evaluable. Of the 101 evaluable rules, 86 refer to a specific group in the main stratum, 7 to protonated amines and are scored in the salt stratum only, and 8 are umbrella rules (``any carbonyl, 1600 to 1900'') that the source uses as a first pass; these are scored against union groups and reported in their own tier, not compared with group-specific rules. Thirty-one rules state an intensity and 13 a shape.

The table was frozen, together with every choice described below, before the confirmatory data were examined, and the frozen version is tagged in the repository (\texttt{rules-frozen-v1}). No window, qualifier or threshold was changed afterwards; where a rule disagrees with the data, the disagreement is reported.

\subsection{Ground truth}
\label{sec:truth}

A spectrum is positive for a group when the depositor's structure contains a substructure that matches a SMARTS pattern for that group. Seventeen patterns were taken from Punjabi et al.~\cite{punjabi2025chemotion}, who adapted them from Fine et al.~\cite{fine2020spectral}; 43 were written for this work so that every evaluable rule has a target that means what the source means. The definitions follow the source's own words as literally as possible: a primary alcohol is a hydroxyl on a CH$_2$; ``saturated'' means that the atom attached to the group is sp$^3$ and ``aromatic'' that it is aromatic, so the two classes are exclusive and neither includes $\alpha,\beta$-unsaturated carbonyls, which count only for the parent group; amines and amides exclude N-N, N-O and acyl nitrogen. The 60 patterns are listed in the Supplementary Information and in the repository, with the number of matches in each collection.

Structures were standardised with RDKit~\cite{rdkit} (largest fragment for the main stratum, charges as deposited) and identified by InChIKey. The ground truth is therefore the structure the depositor recorded; we did not verify structures against the spectra, since that would presuppose the rules under test. Deposition errors exist; their expected effect is to add noise and to lower measured performance, not to favour any rule.

\subsection{Spectral collections}
\label{sec:data}

\paragraph{Chemotion.} The primary collection is the open IR dataset released with the Chemotion repository~\cite{punjabi2025chemotion} (RADAR4Chem, CC BY-SA 4.0), 4,183 JCAMP-DX files from synthesis laboratories, deposited between 2016 and 2024, each linked to the depositor's structure. Nearly all were recorded on Bruker Alpha ATR instruments with a diamond crystal at 4~\wn{} resolution, from 374 to 3997~\wn{}; the collection is an ATR collection and its intensities carry the ATR wavelength dependence. Chemotion exports each measurement twice (raw and peak-picked), so one file was kept per analysis; files without data points, with fewer than 500 points, not covering 650 to 3800~\wn{}, with constant intensity or with an unparsable structure were removed. This left 2,113 spectra of 1,921 distinct structures. Spectra of metal complexes (150) were excluded; spectra of salts and other multi-fragment entries without metals (40) form a salt stratum used only for the amine salt rules, and the remaining 1,923 spectra form the main stratum.

Twenty percent of the distinct structures (seed 20260914) were set aside as a development set on which the pipeline was built and checked; duplicates of one structure never straddle the two sets. All decisions in this section were taken on the development set. The confirmatory set, 1,692 spectra of which 1,539 are in the main stratum, was analysed once, after the freeze, and is the only Chemotion set reported.

\paragraph{NIST.} The replication collection is the infrared section of the NIST Chemistry WebBook (Standard Reference Database 69)~\cite{nist_webbook}, from which every condensed-phase spectrum with data points was retrieved with a script that is provided with the repository (the spectra themselves are subject to NIST's terms of use and are not redistributed). These are transmission spectra, mostly from the Coblentz Society and NIST collections recorded before 1990 on grating or prism instruments, of KBr discs (1,742 of the primary spectra), mulls (1,360), solutions (1,254) and neat liquids or films (1,126). Records were kept if the abscissa was in wavenumber or micrometres (the latter converted), the ordinate in transmittance or absorbance, the state condensed, the range covering 700 to 3800~\wn{} after conversion, and a structure available from the WebBook; 5,894 spectra passed. Where a species had several spectra one was kept, preferring wavenumber over micrometre abscissae (which are digitised and resampled) and, in order, KBr, liquid, solution and mull preparations. The main stratum holds 5,124 spectra of 5,122 structures; 66\% contain an aromatic ring, against 93\% in Chemotion. Because many prism-era spectra begin at 650 to 700~\wn{}, each rule was scored only on the spectra whose measured range covers every window it uses; this rule was applied to both collections and changes nothing in Chemotion.

\subsection{Processing and band detection}
\label{sec:bands}

Spectra were converted to absorbance ($A = -\log_{10} T$, with $T$ floored at $10^{-4}$), interpolated onto a 2~\wn{} grid from 400 to 4000~\wn{}, and baseline-corrected with asymmetric least squares~\cite{eilers2005baseline} (smoothness $10^6$, asymmetry 0.001, ten iterations), which follows the underside of the spectrum without cutting into bands. Each spectrum was then scaled so that its strongest point between 650 and 3800~\wn{} equals 1; every intensity in this work is relative to that point. Bands are the local maxima of the corrected spectrum with a prominence, the height above the higher of the two flanking minima, of at least 0.01; each is described by its apex position, its height relative to the strongest band and its full width at half prominence. Intensity classes follow the three levels of the source: strong, at least 0.6 of the strongest band; medium, at least 0.2; weak, below. Shape classes use the width: sharp below 30~\wn{}, broad from 30 to 100, very broad above.

The prominence threshold at which a candidate counts as a band is the one parameter with a visible effect on the results, and it was fixed on the development set: 0.02 of the strongest band is the primary value, because real diagnostic bands in this collection can be that weak (the C$\equiv$N stretch of Section~\ref{sec:nitrile} is the example), and 0.01 is at the noise floor of some spectra. Every result was also computed at 0.01, 0.05 and 0.10, and the sweep is given in the Supplementary Information.

\subsection{Scoring a rule}
\label{sec:scoring}

A rule is satisfied by a spectrum when a band with prominence at or above the threshold has its apex inside the window and, in the full definition, its intensity class and shape class match the qualifiers the source states (a ``medium'' qualifier accepts strong or medium; a ``broad'' qualifier accepts broad or very broad). Three nested definitions were computed for every rule, position only, position with intensity, and position with intensity and shape; the full definition is primary because it is what the source says, and the position-only result is reported alongside it. For the benzene substitution rules, a required second window must also contain a band and a forbidden window must contain none, at the same threshold. Each rule was scored on the spectra of its stratum whose measured range covers all its windows.

\subsection{Statistics}
\label{sec:metrics}

For each rule and definition a 2$\times$2 table was formed from the rule outcome and the ground truth, and sensitivity, specificity, predictive values and the positive and negative likelihood ratios~\cite{deeks2004likelihood} were computed, with 95\% confidence intervals on the logarithmic scale~\cite{simel1991likelihood} and a continuity correction of 0.5 when a cell was empty. Because spectra of related compounds are not independent (the 1,923 Chemotion spectra fall into 545 Murcko scaffold~\cite{bemis1996murcko} families, the largest holding 193 spectra), a second interval for \LRp{} was obtained by a cluster bootstrap that resamples scaffold families rather than spectra (500 resamples, fixed seed); it is the interval shown in the figures. Rules with fewer than ten positive spectra in a collection are reported but marked underpowered and not interpreted. For each rule with at least five false positives, the groups over-represented among the false positives relative to the true negatives were listed by odds ratio, as the basis of the failure analysis. No correction for multiple comparisons was applied and none is needed: the study estimates a likelihood ratio for each rule and does not test whether any rule is ``significant''. Four hypotheses about the pattern of results were written down before the confirmatory run and are given, with their outcome, in the Supplementary Information.

\subsection{Software and availability}
\label{sec:software}

The pipeline is written in Python with RDKit~\cite{rdkit} for substructure matching and scaffolds, SciPy~\cite{virtanen2020scipy} for baseline correction and peak detection, and the jcamp package for file parsing. The frozen rule table, the SMARTS definitions, all scripts, the list of NIST identifiers used and the per-rule result tables at every threshold and definition are available at \url{https://github.com/Kementsu/ir-correlation-rules-paper}, tagged \texttt{rules-frozen-v1}. The Chemotion spectra are available from RADAR4Chem; the NIST spectra can be re-downloaded with the script provided. A large language model was used to draft text and analysis code under the author's direction, as stated in the declaration at the end of the article.
